% Generated by src/build_paper_probe_tables.py.
% Source: results/tables/gemma_scope_causality_<label>_local.csv (mode==steering, alpha=+0.10), Gemma-3-12B/27B only.
% table* + resizebox: 6 direction columns overflow a single column.
% [tp]: short table, Results layout pass (see step_logs/STEP_LOG.md).
\begin{table*}[tp]
\centering
\small
\setlength{\tabcolsep}{4pt}
\resizebox{\textwidth}{!}{%
\begin{tabular}{@{}llrrrrrr@{}}
\toprule
Model & Axis & Dense target & SAE decoded & Axis-specific & Shared & Opposing axis & Random \\
\midrule
Gemma-3-12B & Warmth & \textbf{+3.88} & +2.30 & +2.10 & +2.16 & +3.24 & -0.32 \\
Gemma-3-12B & Competence & +2.01 & +3.69 & +3.12 & +2.34 & +2.38 & +2.88 \\
Gemma-3-27B & Warmth & +1.02 & -0.19 & -1.38 & +0.69 & +2.86 & +1.34 \\
Gemma-3-27B & Competence & +0.21 & +3.53 & +4.36 & +1.61 & +0.31 & +2.67 \\
\bottomrule
\end{tabular}%
}
\caption{\textbf{Direction specificity at the local endpoint ($\alpha=+0.10$).} Change in the held-out Yes-versus-No logit margin for six directions, Gemma-3-12B and Gemma-3-27B. Dense target is the same mean-difference direction used everywhere else in this study; SAE decoded, axis-specific, and shared are built from a Gemma Scope~2 sparse decomposition; opposing axis steers along the other concept's direction; random is a single orthogonalized control. Bold marks the row's largest effect. Dense target is largest in only 1 of the four model-axis rows; the other three are matched or exceeded by at least one unrelated direction, most often the SAE-decoded or axis-specific component, which does not support a clean specificity reading and is discussed in the text alongside the shared warmth-competence overlap already noted in PCA denoising.}
\label{tab:concept_direction_specificity}
\end{table*}
